\section{Local Zeckendorf Normalization and Uniqueness of the Fibonacci Weighting}\label{app:canonical-fold}

\subsection{Local Normalization Description of the Fold}

This subsection records a local description of the canonical fold from
Definition~\ref{def:canonical-fold}.

Let \(\mathcal{D}\) be the set of finitely supported digit configurations
\[
a=(a_1,a_2,\dots),\qquad a_k\in\mathbb{N},
\]
and define the Fibonacci value
\[
\mathrm{Val}(a):=\sum_{k\ge 1} a_k F_{k+1}.
\]
Embed \(\omega=\omega_1\cdots \omega_m\in\Omega_m\) as
\[
\iota_m(\omega):=(\omega_1,\dots,\omega_m,0,0,\dots)\in\mathcal{D}.
\]
Consider the following value-preserving local rewrites:
\[
e_k+e_{k+1}\to e_{k+2},
\qquad
2e_1\to e_2,
\qquad
2e_2\to e_1+e_3,
\qquad
2e_k\to e_{k-2}+e_{k+1}\ \ (k\ge 3),
\]
where \(e_k\) is the \(k\)-th unit vector in \(\mathcal{D}\).

\begin{proposition}[Order-independent local normalization]\label{prop:canonical-fold-rewrite}
Starting from \(\iota_m(\omega)\), every maximal rewrite sequence terminates at
the same normal form \(\mathrm{NF}(\iota_m(\omega))\in\{0,1\}^{\mathbb{N}}\)
with no adjacent \(1\)'s. This normal form is the Zeckendorf expansion
\(Z(N(\omega))\). Consequently,
\[
\Fold_m(\omega)=\pref_m\bigl(\mathrm{NF}(\iota_m(\omega))\bigr)
\]
is independent of rewrite order.
\end{proposition}

\begin{proof}
Each rewrite preserves \(\mathrm{Val}\) because it encodes one of the Fibonacci
identities
\[
F_{k+1}+F_{k+2}=F_{k+3},
\qquad
2F_2=F_3,
\qquad
2F_3=F_2+F_4,
\qquad
2F_{k+1}=F_{k-1}+F_{k+2}\ \ (k\ge 3).
\]
Termination follows from a lexicographic descent argument. Let
\[
A(a):=\sum_{k\ge 1} a_k,
\qquad
B(a):=\sum_{k\ge 1} k\,a_k,
\qquad
C(a):=a_2.
\]
The rules \(e_k+e_{k+1}\to e_{k+2}\) and \(2e_1\to e_2\) decrease \(A\); the
rules \(2e_k\to e_{k-2}+e_{k+1}\) for \(k\ge 3\) preserve \(A\) and decrease
\(B\); and the rule \(2e_2\to e_1+e_3\) preserves \(A\) and \(B\) but decreases
\(C\). Hence \((A,B,C)\) decreases lexicographically at each step, so rewriting
must terminate.

A normal form has no adjacent \(1\)'s and no digit \(\ge 2\), hence is a
Zeckendorf representation. Since every rewrite preserves \(\mathrm{Val}\) and
Zeckendorf representation is unique \cite{Zeckendorf1972}, all terminating
branches have the same normal form. The final formula for \(\Fold_m\) is then
just the definition of the canonical fold through prefix truncation.
\end{proof}

\subsection{Uniqueness of the Fibonacci Weighting}

The preceding subsection shows that the canonical fold is implemented by
Zeckendorf normalization. A natural question is whether the choice of
Fibonacci weights \(F_{k+1}\) is itself forced. The following theorem shows
that it is, within a natural axiom class.

\begin{theorem}[Contiguous numeration uniqueness]
\label{thm:fibonacci-weight-uniqueness}
Let \(w_2,w_3,\dots\) be positive integers. Define the valuation
\[
V_m^{(w)}(x_1\cdots x_m):=\sum_{k=1}^m x_k\,w_{k+1}.
\]
If for every \(m\ge 1\) the map \(V_m^{(w)}\) restricts to a bijection from
\(X_m\) onto \(\{0,1,\dots,F_{m+2}-1\}\), then \(w_{k+1}=F_{k+1}\) for all
\(k\ge 1\).
\end{theorem}

\begin{proof}
For \(m=1\), \(X_1=\{0,1\}\) and the bijection condition gives
\(\{0,w_2\}=\{0,1\}\), so \(w_2=1=F_2\).

For \(m=2\), \(X_2=\{00,01,10\}\) and \(\{0,w_2,w_3\}=\{0,1,2\}\); since
\(w_2=1\), one obtains \(w_3=2=F_3\).

Proceed by induction. Suppose \(w_2=F_2,\dots,w_m=F_m\) for some
\(m\ge 3\). Decompose \(X_m\) by its last symbols:
\[
X_m
=
\{u0: u\in X_{m-1}\}
\;\sqcup\;
\{v01: v\in X_{m-2}\}.
\]
On the first part, \(V_m^{(w)}(u0)=V_{m-1}^{(w)}(u)\). On the second,
\(V_m^{(w)}(v01)=V_{m-2}^{(w)}(v)+w_{m+1}\). By the inductive hypothesis,
\[
V_{m-1}^{(w)}(X_{m-1})=\{0,\dots,F_{m+1}-1\},
\quad
V_{m-2}^{(w)}(X_{m-2})=\{0,\dots,F_m-1\}.
\]
Hence
\[
V_m^{(w)}(X_m)
=
\{0,\dots,F_{m+1}-1\}
\;\sqcup\;
\{w_{m+1},\dots,w_{m+1}+F_m-1\}.
\]
For this to equal \(\{0,\dots,F_{m+2}-1\}=\{0,\dots,F_{m+1}+F_m-1\}\), the
second interval must start at \(F_{m+1}\), giving \(w_{m+1}=F_{m+1}\).
\end{proof}

\begin{corollary}[Uniqueness of the canonical residue fold]
\label{cor:canonical-fold-uniqueness}
Let \(G_m:\Omega_m\to X_m\) be a deterministic map satisfying:
\begin{enumerate}
\item \(G_m(x)=x\) for all \(x\in X_m\).
\item There exist positive integer weights \(w\) such that
  \(V_m^{(w)}|_{X_m}\) is a bijection onto \(\{0,\dots,F_{m+2}-1\}\).
\item For every \(\omega\in\Omega_m\),
  \(V_m^{(w)}(G_m(\omega))\equiv \sum_{k=1}^m\omega_k w_{k+1}
  \pmod{F_{m+2}}\).
\end{enumerate}
Then \(G_m=\Fold_m\).
\end{corollary}

\begin{proof}
Theorem~\ref{thm:fibonacci-weight-uniqueness} forces \(w_{k+1}=F_{k+1}\).
Condition~2 then provides a unique legal representative in each residue class
modulo \(F_{m+2}\), and condition~3 selects that representative. This is
exactly the residue characterization of
Proposition~\ref{prop:fold-residue-characterization}.
\end{proof}
